\documentclass[crop,tikz]{standalone}
\usepackage[x11names]{xcolor}
\usepackage{tikz}
\usepackage[scaled]{helvet}
\renewcommand*\familydefault{\sfdefault}

\begin{document}

\pgfdeclarelayer{background}
\pgfsetlayers{background,main}

\usetikzlibrary{arrows,shapes,positioning,calc}

\tikzstyle{package} = [
  rectangle, 
  font=\large,
  minimum height=2.0em,
  rounded corners=0.75em,
  ultra thick]

\tikzstyle{sub-pylith} = [
  rectangle, 
  font=\large,
  minimum height=2.0em,
  ultra thick,
  text width=12em, 
  anchor=west]
\tikzstyle{pylith} = [
  package, 
  font=\bf\large,
  minimum width=16em, 
  draw=Red2!80!black,
  top color=Red1!50!white,
  bottom color=Red3]
\tikzstyle{cxx} = [
  draw=Orange2!80!black,
  top color=Orange1!20!white,
  bottom color=Orange2,
  ]
\tikzstyle{swig} = [
  draw=Green2!80!black,
  top color=Green1!50!white,
  bottom color=Green3,
  ]
\tikzstyle{python} = [
  draw=DodgerBlue2!80!black,
  top color=DodgerBlue1!50!white,
  bottom color=DodgerBlue3,
  ]

\tikzstyle{python-package} = [
  package,
  minimum width=8em,
  text depth=4em]


\tikzstyle{arrowto} = [>=latex, ->, very thick, color=black]

\begin{tikzpicture}[node distance=0.5ex and 2.0ex]

  \begin{pgfonlayer}{background}
    
    % Component
    \node (pylith) [pylith, text depth=16em] {{\bfseries PyLith}};
    
  \end{pgfonlayer}
  
  % Python
  \node (python) [sub-pylith, python, below=of pylith.north, yshift=-2em, text depth=4em]
    {{\bfseries Python}\\[\baselineskip] User Input};
  \node (swig) [sub-pylith, swig, below=of python] {{\bfseries SWIG}};
  \node (cxx) [sub-pylith, cxx, below=of swig]
    {{\bfseries C/C++}\\[\baselineskip] Computational Engine};


% Dependencies
\node (pyre) [package, draw=DodgerBlue2!80!black, top color=DodgerBlue2, bottom color=Orange2, xshift=2em, right=of python.north east, anchor=north west] {Pyre};
\node (h5py) [python, package, right=of pyre] {h5py};
\node (numpy) [python, package, right=of h5py] {numpy};

\node (spatialdata) [package, draw=Orange2!80!black, top color=DodgerBlue2, bottom color=Orange2, xshift=2em, right=of cxx.south east, anchor=south west] {SpatialData};
\node (petsc) [cxx, package, right=of spatialdata] {PETSc};
\node (proj) [cxx, package, right=of petsc] {Proj};
\node (netcdf) [cxx, package, right=of proj] {NetCDF};
\node (hdf5) [cxx, package, right=of netcdf] {HDF5};
\node (mpi) [cxx, package, right=of hdf5] {MPI};

  % Dependency
  \draw[arrowto] (python.east) to (spatialdata.north);
  \draw[arrowto] (python.east) to[bend right] (pyre.south);
  \draw[arrowto] (python.east) to[bend right] (h5py.south);
  \draw[arrowto] (python.east) to[bend right] (numpy.south);
  \draw[arrowto] (pyre.south) to[bend right] (h5py.south);
  \draw[arrowto] (pyre.south) to (mpi.north);
  \draw[arrowto] (h5py.south) to (hdf5.north);
  \draw[arrowto] (h5py.south) to[bend right] (numpy.south);

  \draw[arrowto] (cxx.east) to[bend left] (petsc.north);
  \draw[arrowto] (cxx.east) to[bend left] (netcdf.north);
  \draw[arrowto] (cxx.east) to[bend left] (hdf5.north);
  \draw[arrowto] (cxx.east) to[bend left] (mpi.north);
  \draw[arrowto] (cxx.east) to[bend left] (spatialdata.north);
  \draw[arrowto] (spatialdata.north) to (pyre.south);
  \draw[arrowto] (spatialdata.north) to[bend left] (proj.north);
  \draw[arrowto] (petsc.north) to[bend left] (hdf5.north);
  \draw[arrowto] (petsc.north) to[bend left] (mpi.north);  
  \draw[arrowto] (netcdf.north) to[bend left] (hdf5.north);
  \draw[arrowto] (hdf5.north) to[bend left] (mpi.north);

\end{tikzpicture}

\end{document}
